The \textbf{CubeSat Camera System (C2S)} is the proposed mission payload for the Perth Aerospace Student Team's (PAST) first CubeSat. This inaugural mission represents a critical milestone in the team's broader goal of developing low-cost, modular spacecraft capable of performing meaningful scientific and engineering demonstrations in Low Earth Orbit (LEO). C2S is designed to operate as an Earth observation payload, taking photos of land- and ocean-scapes. Beyond its imaging objectives, the system will serve as a testbed for evaluating the viability of student-developed space hardware in a real orbital environment. The payload aims to balance performance, reliability, and resource constraints typical of CubeSat-class missions.
\nl
This document outlines the key design drivers, system architecture, component selection strategies, and trade-offs that inform the development of the CubeSat Camera System prototype. It also provides an overview of the integration challenges and planned validation strategies for the final model.
\nl
The project's development is encapsulated through three versions. Versions 0 and 0.1 focused on developing a prototype to confirm and validate functionality, whereas Version 1 regarded the flight prototype. Versions 0 \& 1 were uniquely different in their system requirements - the former having very little hardware restrictions with more features tailored to development like scroll wheels and visual displays, and the latter optimising performance and reducing size. 
\nl
Where appropriate, this document often refers to prerequisite knowledge which can be found in the \autoref{sec:appendices}. Please refer to the relevant appendix to fill gaps in understanding.